% Part: second-order-logic
% Chapter: sol-and-set-theory
% Section: cardinalities

\documentclass[../../../include/open-logic-section]{subfiles}

\begin{document}

\olfileid{sol}{set}{crd}

\olsection{أعداد المجموعات الأصلية}

\begin{explain}
مثلما نستطيع التعبير عن كون المجال منتهيًا أو لانهائيًا،
!!{enumerable} أو~!!{nonenumerable}، نستطيع تعريف خاصية كون مجموعة
جزئية من~$\Domain{M}$ منتهية أو لانهائية، !!{enumerable} أو
!!{nonenumerable}.
\end{explain}

\begin{prop}
الصيغة~$\fn{Inf}(X) \ident$
\begin{multline*}
\lexists[u][(\lforall[x][(X(x) \lif X(u(x)))] \land {}\\
  \lforall[x][\lforall[y][(\eq[u(x)][u(y)] \lif 
      \eq[x][y])]] \land {} \\\lexists[y][(X(y) \land \lforall[x][(X(x)
      \lif \eq/[y][u(x)])])])]
\end{multline*}
تُرضى بالنسبة إلى تعيين متغيرات~$s$ إذا وفقط إذا كانت~$s(X)$
لانهائية.
\end{prop}

\begin{prop}
الصيغة~$\fn{Count}(X) \ident $
\begin{multline*}
(\lnot\lexists[x][X(x)] \lor {}\\
\lexists[z][\lexists[u][(X(z) \land 
    \lforall[x][(X(x) \lif X(u(x)))] \land {} \\ \lforall[Y][((Y(z) \land
      \lforall[x][(Y(x) \lif Y(u(x)))]) \lif X \subseteq Y)])]])
\end{multline*}
تُرضى بالنسبة إلى تعيين متغيرات~$s$ إذا وفقط إذا كانت~$s(X)$
!!{enumerable}.
\end{prop}

نعلم من مبرهنة كانتور أن ثمة مجموعات~!!{nonenumerable}، وأن ثمة، في
الواقع، عددًا لانهائيًا من المستويات المختلفة للأحجام اللانهائية. وتطور
نظرية المجموعات حسابًا كاملًا لأحجام المجموعات، وتسند أعدادًا أصلية
لانهائية إلى المجموعات. وتؤدي الأعداد الطبيعية دور الأعداد الأصلية التي
تقيس أحجام المجموعات المنتهية. أما أصلية المجموعات~!!{denumerable} فهي
الأصلية اللانهائية الأولى، وتسمى~$\aleph_0$ («ألف-صفر»). والحجم
اللانهائي التالي هو~$\aleph_1$، وهو أصغر حجم يمكن أن تكون عليه مجموعة
من دون أن تكون قابلة للتعداد (أي من الحجم~$\aleph_0$). ويمكننا تعريف
«$X$ من الحجم~$\aleph_0$» بـ
$\fn{Aleph}_0(X) \liff \fn{Inf}(X) \land \fn{Count}(X)$. وتكون~$X$
من الحجم~$\aleph_1$ إذا كانت لانهائية، وكانت جميع مجموعاتها الجزئية
منتهية أو من الحجم~$\aleph_0$ أو مكافئة لها عدديًا، ولم تكن هي نفسها
من الحجم~$\aleph_0$. ومن ثم نستطيع التعبير عن ذلك بالصيغة
$\fn{Aleph_1}(X) \ident \fn{Inf}(X) \land
\lforall[Y][(Y \subseteq X \lif
  (\lnot\fn{Inf}(Y) \lor \fn{Aleph}_0(Y) \lor \cardeq{X}{Y}))] \land \lnot
\fn{Aleph}_0(X)$. ويعرّف كون المجموعة من الحجم~$\aleph_2$ على نحو
مماثل، وهكذا.

وثمة حجم له أهمية خاصة، هو ما يسمى أصلية المتصل. وهو حجم
$\Pow{\Nat}$، أو، على نحو مكافئ، حجم~$\Real$. ويمكن أيضًا التعبير في
منطق الرتبة الثانية عن كون مجموعة ما من حجم المتصل، لكن ذلك يتطلب قدرًا
أكبر من العمل.

\end{document}
